% Frontier addendum, 2026-09-04.
% Intended for later integration into arithmetic_cat_map_moments.tex after full manuscript audit.

\section{An exact-baseline density theorem from the global resultant budget}
\label{sec:exact-baseline-density}

We retain the notation of the Cayley-energy reduction.  Let $K$ be an odd
prime, let $p\equiv1\pmod K$ be prime, let $H\leq\mathbb F_p^\times$ be the
subgroup of order $K$, and put
\[
 A=c(H\setminus\{1\}),\qquad c(h)=\frac{h+1}{h-1}.
\]
For nonzero differences define
\[
 E_K^{\rm nz}(p)=\sum_{d\ne0}r_{A-A}(d)^2,
 \qquad
 B_K=(K-1)(2K-5),
\]
and let
\[
 \Delta_K(p)=E_K^{\rm nz}(p)-B_K\in\mathbb Z_{\ge0}.
\]
The global resultant argument established elsewhere in this manuscript gives
\begin{equation}
 \sum_{p\equiv1\,({\rm mod}\,K)}\Delta_K(p)\log p
 \le N_K C_K,
 \label{eq:global-resultant-budget-addendum}
\end{equation}
where
\[
 N_K=K^3-5K^2+6K+1,
 \qquad
 C_K=(K-1)\log 12-3\log K.
\]
In particular $N_K<K^3$ and $C_K<K\log12$ for $K\ge3$.

Fix $0<\theta<1/2$, write $Q=X^\theta$, and define the dyadic split family
\[
 \mathcal D_\theta(X)=\{(K,p):K,p\ {\rm prime},\ Q/2<K\le Q,
 X/2<p\le X,\ p\equiv1\pmod K\}.
\]
Let
\[
 \mathcal B_\theta(X)=\{(K,p)\in\mathcal D_\theta(X):\Delta_K(p)>0\}.
\]
The Bombieri--Vinogradov theorem together with the prime number theorem gives
\begin{equation}
 |\mathcal D_\theta(X)|\sim
 \frac{\log2}{2\theta}\frac{X}{(\log X)^2}.
 \label{eq:dyadic-design-count-addendum}
\end{equation}

\begin{theorem}[density-one exact Cayley baseline]
\label{thm:exact-baseline-density}
For every fixed $0<\theta<1/2$,
\[
 |\mathcal B_\theta(X)|
 \ll_\theta \frac{X^{5\theta}}{(\log X)^2},
 \qquad
 \frac{|\mathcal B_\theta(X)|}{|\mathcal D_\theta(X)|}
 \ll_\theta X^{5\theta-1}.
\]
Consequently, for every fixed $0<\theta<1/5$, asymptotically almost every
$(K,p)\in\mathcal D_\theta(X)$ satisfies
\[
 \Delta_K(p)=0
\]
exactly, and hence
\[
 E_+(A)=3(K-1)(K-2).
\]
\end{theorem}

\begin{proof}
Fix a prime $K\in(Q/2,Q]$.  If $p\in(X/2,X]$ is bad, then the integrality and
nonnegativity of $\Delta_K(p)$ imply $\Delta_K(p)\ge1$.  Thus every bad prime
contributes at least $\log(X/2)$ to the left-hand side of
\eqref{eq:global-resultant-budget-addendum}.  Therefore
\[
 \#\{p\in(X/2,X]:p\equiv1\pmod K,\ \Delta_K(p)>0\}
 \le \frac{N_KC_K}{\log(X/2)}
 \ll \frac{K^4}{\log X}.
\]
Summing over prime $K\in(Q/2,Q]$ and applying the prime number theorem with
partial summation,
\[
 |\mathcal B_\theta(X)|
 \ll \frac1{\log X}
       \sum_{\substack{Q/2<K\le Q\\K\ {
m prime}}}K^4
 \ll \frac{Q^5}{\log Q\log X}
 \ll_\theta \frac{X^{5\theta}}{(\log X)^2}.
\]
Division by \eqref{eq:dyadic-design-count-addendum} yields the second bound.
If $\theta<1/5$, its right-hand side tends to zero.  Finally the zero-difference
class contributes $(K-1)^2$ and the forced nonzero class contributes
$(K-1)(2K-5)$, giving
\[
 E_+(A)=(K-1)^2+(K-1)(2K-5)=3(K-1)(K-2)
\]
on the good subfamily.
\end{proof}

The theorem is deliberately an additive-energy statement, not a microscopic
Bohigas--Giannoni--Schmit theorem.  It strengthens the previously obtained
almost-everywhere $o(K^2)$ excess to exact vanishing on the smaller range
$\theta<1/5$.  The exponent $1/5$ reflects the current $O(K^4)$ total
logarithmic resultant budget.  A constant-size collision symmetry orbit cannot
change this exponent; an improvement requires either a smaller global
resultant budget, a growing compulsory multiplicity for a bad reduction, or a
new restriction on the split prime divisors of the normalized resultants.

For the ordinary full-annihilator normalized Jacobi trace, the exact-baseline
identity also gives
\[
 \kappa(K,p)
 =\frac{(p-1)\bigl(3p(K-2)-(K-1)^3\bigr)}
        {(K-1)(p-K+1)^2}
 =3+O\!\left(K^{-1}+K^2/p\right).
\]
This fourth-moment ratio does not imply a central limit theorem and is not a
local eigenphase statistic.
